\section{Metodología Experimental}

Se ensambló un transformador utilizando un núcleo de hierro acerado y un juego de bobinas de 300 y 600 espiras. El sistema se conectó a una fuente de voltaje alterno variable $(0 - 250\,\mathrm{V})$, empleando multímetros digitales para la medición directa de las variables eléctricas.

\subsection{Caracterización en Circuito Abierto (Relación de Transformación)}

\begin{enumerate}
    \item \textbf{Configuración:} Se estableció el circuito con el secundario abierto $(R=\infty)$, utilizando inicialmente una relación de vueltas de $N_p = 600$ y $N_s = 300$.
    
    \item \textbf{Toma de datos:} Se varió el voltaje del primario $(V_p)$ de forma gradual. Para cada incremento, se midieron directamente el voltaje inducido en el secundario $(V_s)$ y la corriente de vacío en el primario $(I_p)$.
    
    \item \textbf{Repetición:} El procedimiento se repitió modificando el número de espiras en el secundario para observar el cambio en la constante de transformación:
    
    \begin{equation}
        K=\frac{N_s}{N_p}
    \end{equation}
\end{enumerate}
\begin{figure}[H]
    \centering
    \includegraphics[width=0.45\linewidth]{montaje.jpeg}
    \caption{Montaje Experimental}
    \label{fig:Montaje}
\end{figure}
\subsection{Ensayo con Carga (Análisis de Eficiencia)}

\begin{enumerate}
    \item \textbf{Conexión de carga:} Se integró un reóstato de $1440\,\Omega$ en el circuito del secundario para actuar como carga resistiva variable.
    
    \item \textbf{Estabilización del voltaje:} Se fijó un voltaje constante en el devanado primario.
    
    \item \textbf{Medición de potencia:} Al variar la resistencia del reóstato, se registraron simultáneamente los voltajes eficaces $(V_p, V_s)$ y las corrientes $(I_p, I_s)$. A partir de estas mediciones se calcularon indirectamente la potencia de entrada $(P_p)$, la potencia de salida $(P_s)$ y la eficiencia energética $(\eta)$ del transformador.
\end{enumerate}

\subsection{Análisis de Datos y Limitaciones}

El análisis experimental se fundamentó en la aplicación de técnicas de regresión lineal para verificar la proporcionalidad entre los voltajes del devanado primario y secundario, así como para evaluar la transferencia de potencia en el transformador. Las principales fuentes de incertidumbre identificadas correspondieron a la resistencia interna de los conductores, las pérdidas propias del núcleo magnético y la precisión instrumental de los multímetros utilizados. Estas contribuciones fueron consideradas mediante la Teoría de Propagación de errores con el fin de estimar la confiabilidad de los resultados obtenidos.
